\documentclass[11pt]{article}
\usepackage[margin=1in]{geometry}
\usepackage{amsmath,amssymb,amsthm,hyperref}
\newtheorem{theorem}{Theorem}
\newtheorem{corollary}[theorem]{Corollary}
\begin{document}
\setlength{\emergencystretch}{2em}
\title{Ordinary binomial SFC at equal prime exponent and gap}
\author{}\date{}\maketitle
An exact mixed-rank identity is available for every shape offset
$1\le h\le d$. When the lower exponent and gap are the same odd prime,
the correction term disappears in that characteristic. A cyclic
Krylov matrix then excludes every consecutive common root.
The prime two case follows from the retained lower-two theorem.
All claims about SFC below use the ordinary factorial functional
$\mathcal L_1(X^j)=j!$; no unrestricted Gamma-shape conclusion is made.

For $n\ge1$, $d\ge2$, and $\gamma>0$, write
\[
 T_{n,\gamma,d}(t)=\sum_{k=0}^n\binom nk(\gamma)_{dk}t^k,\quad
 \epsilon=(-1)^{d+1},\quad
 R=z^{dn}T_{n,\gamma,d}(\epsilon z^{-d}),\quad
 M=\gamma+dn,\quad E=M-z\partial_z.
\]
Let $B_j=B_{j,\gamma,d}$, always at the SAME $\gamma$, where
\[
 B_0=1,\qquad B_{j+1}=(z+dj+\gamma+d-1)B_j-zB_j',
\]
and put
\[
 L_j(z)=z^j(\partial_z-1)^jB_{n+1}(z).
\]
The cyclic column of a polynomial means its remainder coefficient
column modulo $z^d-u$, in basis $1,z,\ldots,z^{d-1}$.

\begin{theorem}[A mixed-rank criterion for offsets up to the gap]
\label{thm:offset-rank}
Fix an integer $1\le h\le d$. Put $C=M+d$, $A_h=(C)_h$ and
\[
 V_h(E)=\frac{(E)_h-(C)_h}{E-C}
       =\sum_{k=0}^{h-1}v_k(E)_k,\qquad
 \kappa_h=\frac{d(n+1)}{A_h}.
\]
The quotient is a polynomial in $\mathbb Z[C,E]$, with exact coefficients
\[
 v_k=(C+k+1)_{h-1-k}\in\mathbb Z_{\ge0}[C]\quad(0\le k<h).
\]
Set $v_k=0$ for $k\ge h$.
For $u\ne0$, the two equalities
\[
 T_{n,\gamma,d}(\epsilon/u)
  =T_{n+1,\gamma+h,d}(\epsilon/u)=0
\]
hold if and only if the $d$ by $(d-1)$ cyclic matrix whose columns are
\[
 J_j=L_j-\epsilon\kappa_h u\,v_{d-1-j}B_n,
                \qquad 0\le j\le d-2,                         \tag{1}
\]
has rank less than $d-1$. In (1), $u$ is a scalar in the cyclic fiber.
\end{theorem}
\begin{proof}
The coefficients follow from a finite telescoping identity.
Set $Q_k=(E)_k(C+k)_{h-k}$ for $0\le k\le h$.
Then
\[
 Q_{k+1}-Q_k=(E-C)(E)_k(C+k+1)_{h-1-k}.
\]
Summing from $k=0$ to $h-1$ proves the displayed expansion of $V_h$
and its coefficient assertions.
Set $K=(E)_dR$. The lower Euler equation and commutator give
\[
 (E-M)K=-\epsilon z^d(E-\gamma)R,\qquad
 Ez^d=z^d(E-d).
\]
Let $F(E)=(E+d)_h$. Polynomial division gives
$F(E)-F(M)=(E-M)G(E)$, where $G(E-d)=V_h(E)$.
Thus the exact reduction is
\[
 F(E)K=A_hK-\epsilon z^dV_h(E)(E-\gamma)R.
\]
Pascal's identity and Gamma shifting give the upper moment operator
\[
 (\gamma)_h z^{d(n+1)}T_{n+1,\gamma+h,d}(\epsilon z^{-d})
       =[z^d(E)_h+\epsilon(E)_{d+h}]R.
\]
Using $(E)_{d+h}=F(E)(E)_d$ and
$(E)_h=A_h+(E-C)V_h(E)$ reduces this to
\[
 \epsilon A_hK+
 z^d[A_h-d(n+1)V_h(E)]R.                                  \tag{2}
\]
At $R=0$, the upper moment therefore vanishes exactly when
\[
 K=\epsilon\kappa_h u V_h(E)R.                             \tag{3}
\]

We use the full rising-jet concomitant identity from the proof of
\texttt{thm:mixed}, equation (3), in the
\href{LOWER_TWO_PRIME_GAP_SFC.tex}{lower-two prime-gap note}.
For each $\xi^d=1$, after removal of a nonzero row gauge and
restriction to $R=0$, that identity is
\[
 B_n(\xi z)(E)_dR-
 \sum_{j=0}^{d-2}L_j(\xi z)(E)_{d-1-j}R=0.                 \tag{4}
\]
The same proof gives full rank $d$ of the concomitant on the full
$(d+1)$-jet, with kernel the actual nonzero jet of $R$.
The restrictions $R=0$ and (3) are independent because
$\deg V_h\le d-1$. Substitution in (4) gives exactly the columns
(1), up to a common column sign and ordering. Hence restricted
rank deficiency is equivalent to the actual jet satisfying both
restrictions. Fourier evaluation at $\xi z$ converts this to the
cyclic-column statement in characteristic zero.
This uses no lower-root simplicity assumption.

We also retain the constant triangularization
\[
 L_{j+1}=z(\partial_z-1)L_j-jL_j,\qquad
 z(\partial_z-1)B_{n+1+r}
  =(M+2d-1+dr)B_{n+1+r}-B_{n+2+r}.                         \tag{5}
\]
It shows that $L_0,\ldots,L_{d-2}$ are a constant triangular
combination of $B_{n+1},\ldots,B_{n+d-1}$ with diagonal
$1,-1,\ldots,(-1)^{d-2}$.
\end{proof}

\begin{theorem}[Equal prime lower exponent and gap]\label{thm:diagonal-prime}
For every prime $p$ and every nonzero complex binomial
\[
 f(X)=X^p(A+BX^p),
\]
no two consecutive positive ordinary factorial moments vanish.
Equivalently, for every $n\ge1$,
\[
 \gcd\bigl(T_{n,pn+1,p},T_{n+1,p(n+1)+1,p}\bigr)=1
\]
over $\mathbb Q[t]$ and over $\mathbb C[t]$.
\end{theorem}
\begin{proof}
For $p=2$ this is the case $\beta=1$ of
\texttt{thm:prime} in the
\href{LOWER_TWO_PRIME_GAP_SFC.tex}{lower-two prime-gap note}.
Suppose $p$ is odd, and apply Theorem~\ref{thm:offset-rank}
with $d=h=p$, $\gamma=pn+1$, and $M=2pn+1$.
In a complete algebraically closed extension of $\mathbb Q_p$,
every coefficient at degree $k$ in the lower $T$ has valuation
at least $k$: each block of $p$ consecutive integer factors
of $(pn+1)_{pk}$ contains a multiple of $p$.
A common root $t$ therefore gives $u=1/t$ with $v_p(u)\ge1$,
by the unique least-term argument for
$u^nT_{n,pn+1,p}(u^{-1})$.

Here $C=2pn+p+1$. The last factor of $(C)_p$ is
$C+p-1=2p(n+1)$, and the other $p-1$ factors are units.
Consequently
\[
 \kappa_p=\frac{p(n+1)}{(C)_p}
     =\frac1{2\prod_{j=0}^{p-2}(C+j)}
\]
is a $p$-adic unit. Each $v_k$ in Theorem~\ref{thm:offset-rank}
is integral, so the perturbations in (1) vanish modulo the
maximal ideal. The reduced cyclic space is $k[z]/(z^p)$.
Its auxiliary operator is
\[
 {\mathsf A}=z-z\partial_z,
 \qquad {\mathsf A}^p={\mathsf A}.
\]
The identity follows, for example, from its cyclic bidiagonal
characteristic polynomial
$\prod_{j=0}^{p-1}(x+j)=x^p-x$.
The seed $1$ is cyclic because its first $p$ powers are monic
of successive degrees. Thus the $p$ distinct eigenvalues are
the elements of $\mathbb F_p$, and $1$ has a nonzero component
in every eigenspace. The nonzero eigenspaces and a Vandermonde
minor show that
\[
 {\mathsf A}^{n+1}1,\ldots,{\mathsf A}^{n+p-1}1
\]
are independent for $n\ge1$.

The triangularization (5) is integral with unit diagonal.
The reduced $J_j$ therefore have rank $p-1$, contradicting
the necessary rank deficiency. Fourier inversion is used only
in characteristic zero; the reduction is directly in cyclic
coefficient columns. The polynomials have rational coefficients;
absence of a common root over this algebraically closed
$p$-adic field implies a nonzero rational resultant, hence
coprimality over $\mathbb C$.

For $AB\ne0$ the exact normalization is
\[
 \mathcal L_1(f^j)=A^j(pj)!\,T_{j,pj+1,p}(B/A).
\]
If one coefficient vanishes, $f$ is a nonzero monomial and all
its positive moments are nonzero. This proves the phase assertion.
\end{proof}

\begin{corollary}[The cubic-gap support three/six]\label{cor:cubic-support}
For every nonzero complex $X^3(A+BX^3)$, no two consecutive
positive ordinary factorial moments vanish.
More explicitly, for arbitrary $\gamma>0$, $n\ge1$, and $u\ne0$,
a common root at $t=1/u$ of $T_{n,\gamma,3}$ and
$T_{n+1,\gamma+3,3}$ is equivalent to rank less than two of
the cyclic columns
\[
 \left[B_{n+1}-\frac{3(n+1)u}{(M+3)(M+4)(M+5)}B_n,\ B_{n+2}\right],
 \qquad M=\gamma+3n .
\]
\end{corollary}
\begin{proof}
The ordinary assertion is Theorem~\ref{thm:diagonal-prime} at $p=3$.
For the explicit criterion, polynomial division in the rising basis gives
\[
 V_3(E)=(E)_2+(M+5)E+(M+4)(M+5).
\]
Also $L_0=B_{n+1}$ and
$L_1=(M+5)B_{n+1}-B_{n+2}$.
Subtracting $(M+5)$ times the first column of (1) from its
second and changing a column sign yields the stated pair.
In the ordinary case $\gamma=3n+1$, its scalar coefficient simplifies to
\[
 \kappa_3=\frac1{4(3n+2)(6n+5)},
\]
a three-adic unit, consistently with the general proof.
\end{proof}

The replay checks exact quotient/remainder identities, their
rising-basis coefficients, both low-offset normalizations, reduced
Krylov ranks, scalar unit cancellations, and bounded literal moment
gcds. The all-offset and all-index statements follow from the
written polynomial and rank proofs. No all-index theorem is inferred
from the finite controls. Unrestricted binomial and arbitrary-support
SFC remain outside this result.
\end{document}
